\documentclass[11pt,fleqn]{exam}
\usepackage[utf8]{inputenc}
\usepackage[T1]{fontenc}
\usepackage{fancyvrb}
\usepackage{amsmath}
\usepackage{amssymb}
\usepackage{hyperref}
\usepackage{algpseudocode}
\usepackage{algorithm,algorithmicx}
\usepackage{comment}
\usepackage{enumitem}
\usepackage[margin=0.75in]{geometry}
\usepackage{tikz}
\usetikzlibrary{arrows,arrows.meta,positioning,intersections,shapes.gates.logic.US,calc}
\algnewcommand\algorithmicforeach{\textbf{for each}}
\algdef{S}[FOR]{ForEach}[1]{\algorithmicforeach\ #1\ \algorithmicdo}
\newcommand{\fillinMCmath}[1]{\begin{tikzpicture}\draw circle [radius=0.5em];\end{tikzpicture}\ #1}
\newcommand{\fillinMCmathsoln}[1]{\begin{tikzpicture}\draw[black, fill=blue] circle [radius=0.5em];\end{tikzpicture}\ #1}
\newcommand{\fillinblank}[1]{\fillinblankmath{\mbox{#1}}}
\newcommand{\fillinblankmath}[1]{\begingroup\setlength{\fboxsep}{1em}\setlength{\fboxrule}
{2pt}\fbox{\LARGE\phantom{#1}}\endgroup}
\newcommand{\fillinblankmathsoln}[1]{\begingroup\setlength{\fboxsep}{1em}\setlength{\fboxrule}
{2pt}\fbox{#1}\endgroup}
\newcommand{\ptsamt}[1]{[#1~points]}
\newcommand{\ptamt}[1]{[#1~point]}

%%% Adding Colour to Questions and Answers
\usepackage{color}

\definecolor{solnblue}{rgb}{0,0,1}
\newenvironment{soln}{\color{solnblue}}{}

%fancyvrb commands
\def\lar{$\leftarrow$}
\def\lesseq{$\leq$}
\def\greq{$\geq$}
\def\alp{$\alpha$}
\def\inf{$\infty$}
\def\ne{$\neq$}

\newcommand{\subsum}{\mbox{SS}}
\newcommand{\sub}{\mbox{Sub}}

%Questions

% Answers
\definecolor{blu}{rgb}{0,0,0.5}
\def\blu#1{{\color{blu}#1}}
\definecolor{gre}{rgb}{0,.3,0}
\def\gre#1{{\color{gre}#1}}
\definecolor{red}{rgb}{0.5,0.0,0}
\def\red#1{{\color{red}#1}}
\def\norm#1{\|#1\|}
%%% End for Colours

\bracketedpoints


\makeatletter
\renewenvironment{solution}{\leavevmode\par\begin{soln}\noindent Solution:}{\end{soln}}
\makeatother
 \newif\ifsolutions\solutionstrue
 
\author{}
\date{}


\title{CPSC 320 2024W1: Assignment 4}

\begin{document}

	\maketitle
	
This assignment is due \textbf{Friday, November 15 at 7 PM}. Late submissions will not be accepted. All the submission and formatting rules for Assignment~1 apply to this assignment as well.  

%------------------------------------------------------------------------------------
\section{List of names of group members (as listed on Canvas)}

Provide the list here. This is worth 1 mark. Include student numbers as a secondary failsafe if you wish.

\section{Statement on collaboration and use of resources}
To develop good practices in doing homeworks, citing resources and acknowledging input from others, please complete the following. This question is worth 2 marks.

\begin{enumerate}
\item All group members have read and followed the guidelines for groupwork on assignments given in the Syllabus).

\fillinMCmath{Yes} \hspace{.5in} \fillinMCmath{No}

\item We used the following resources (list books, online sources, etc. that you consulted):
\item One or more of us consulted with course staff during office hours.

\fillinMCmath{Yes} \hspace{.5in} \fillinMCmath{No}

\item One or more of us collaborated with other CPSC 320 students; none of us took
      written notes during our consultations and we took at least a half-hour break afterwards.

\fillinMCmath{Yes} \hspace{.5in} \fillinMCmath{No}

      If yes, please list their name(s) here:

\item One or more of us collaborated with or consulted others outside of CPSC 320; none of us took written notes during our consultations and we took at least a half-hour break afterwards.

\fillinMCmath{Yes} \hspace{.5in} \fillinMCmath{No}

      If yes, please list their name(s) here:

\end{enumerate}
\newpage

%==========================================
	
	\section{Mystery QuickSelect}
The following variant of the QuickSelect algorithm carefully chooses a pivot, so as to
guarantee that the sizes of Lesser and Greater are at most $7n/10$, and may equal $7n/10$ in the worst case.
It does this by choosing $\lceil n/5 \rceil$ elements in $\Theta(n)$ time (line 7 below), and setting the
pivot to be the median of these elements via a recursive call (line 8).
Exactly how the $\lceil n/5 \rceil$ elements are chosen is not
important.  The rest of the algorithm is identical to \textsc{QuickSelect}.

\vspace{.1in}

\begin{algorithmic}[1]
\Function{MysteryQuickSelect}{$A[1..n], k$}
\State $\triangleright$ return the $k$th smallest element of $A$, where $1\le k \le n$ and all elements of $A$ are distinct
\If{$n==1$}
   \State \Return $A[1]$
\Else
   \State $\triangleright$ the next two lines select the pivot element $p$
   \State create array $A'$ of $\lceil n/5 \rceil$ "carefully choosen" elements of $A$ \Comment this takes $\Theta(n)$ time
    \State $p \gets$ \textsc{MysteryQuickSelect}($A', \lceil |A'|/2 \rceil$) \Comment $A'$ has  $\lceil n/5\rceil$ elements
    \State Lesser  $\gets$ all elements from $A$ less than $p$  \Comment Lesser has $7n/10$ elements in the worst case
    \State Greater $\gets$ all elements from $A$ greater than $p$ \Comment Greater has $7n/10$ elements in the worst case

    \If{|Lessor| $\ge k$}
       \State \Return \textsc{MysteryQuickSelect}(Lesser, $k$)
    \ElsIf{|Lesser| $ = k - 1$}
        \State \Return $p$
   \Else  \hspace{.2in} $\triangleright$ |Lesser| $< k-1$
       \State \Return \textsc{MysteryQuickSelect}(Greater, $k - |\mbox{Lesser}| - 1$)
   \EndIf
\EndIf
\EndFunction
\end{algorithmic}

\begin{questions}
\question[3]
Let $T'(n)$ be the {\em worst-case} run time of \textsc{MysteryQuickSelect}. Complete the following
recurrence for $T'(n)$. by replacing the parts with ???s. You can ignore floors and ceilings.  No justification needed.
 \[
   T'(n) = \left\{ \begin{array}{ll}
               c, & \mbox{ if }  \mbox{ ???} \\
            \mbox{ ???} + cn , & \mbox{ if } n > 1
             \end{array} \right.
   \]


\question[3]
Draw the first three levels (0,1, and 2) of the recursion tree for your recurrence. Label each node with the size of the subproblem it represents, as well as the work done at that node not counting recursive calls. Finally, put the total work per level in a column on the right hand side of your tree. You can provide a hand-drawn figure as long as it is clear and legible. 

\question[2]
What is the worst-case runtime of \textsc{MysteryQuickSelect}? Provide a short justification of your answer.

\question[2]
For comparison, what is the worst case runtime of the {\em original}  \textsc{QuickSelect} algorithm from the worksheet, in which lines 7 and 8 of the pseudocode above are replaced by setting the pivot $p$ to $A[1]$?  Provide a short justification of your answer.

\end{questions}

\newpage
	        
\section{Subway Sub-array}

[Thanks to Peter Gu for contributing this problem.]

\vspace{.1in}

\noindent
Peter has just purchased a sandwich from the Life Building Subway. The sandwich can be represented as an array $A$ of $n$ real-valued quantities. Peter is peculiar, he enjoys partitioning his sandwich into contiguous sub-arrays such that every element from the original array belongs to a sub-array. He then scores each sub-array $A[l,r], 1 \le l \le r \le n$ with the following formula:
\[
    \subsum(l,r) = ax^2 + bx + c,
\]
where $a, b, c \in \mathbb{Z}$ and $x$ is the sum of all elements in the sub-array $A[l..r]$. He would like you to develop an algorithm to find the score of a partition that maximizes the sum of the scores of its subarrays.

More formally, an instance of the problem consists of the array $A[1..n]$ plus the quantities $a, b$, and $c$, and the goal is to determine $\sub(n)$, defined as follows when $n \ge 1$:
\[
\sub(n) = \max_{1\le k \le n}  \;\; \max_{0 = p_0 < p_1 < p_2 < \ldots < p_k = n} \;\; \subsum(p_0+1,p_1) + \subsum(p_1+1,p_2) +  \subsum(p_2+1,p_3) + \ldots \subsum(p_{k-1}+1,p_k).
\]
(Here, $k$ is the number of subarrays in the partition, and for $1 \le i \le k$, the $i$th subarray ranges from $p_{i-1}+1$ to $p_i$.) We'll also define $\sub(0)$ to be 0.
\vspace{.1in}

\noindent
{\bf Example:} Suppose that the array is $[1, 2, -3, 2, 1]$ and $a = 1, b = 1, c = 5$. One possible partition into subarrays is $[1, 2], [-3], [2, 1]$. The respective scores for this partition are $[3^2 + 3 + 5], [3^2 - 3 + 5], [3^2 + 3 + 5]$ and the total sum is $45$. 

Note: No justification is needed for the first five parts of this problem.

\begin{questions}
\question[1]
For the example array above, give an optimal partition and write down its value $\sub(5)$.

\question[3]
Provide pseudocode to calculate all of the quantities $\subsum(l,r)$, for $1\le l \le r \le n$. Your pseudocode should run in $O(n^2)$ time.

\question[3]
Let $n \ge 1$.
Suppose that in an optimal partition, the rightmost subarray is $A[i+1..n]$, where $0 \le i \le n-1$. Write down an expression for the value of $\sub(n)$  in terms of the quantity $\subsum(i+1,n)$ and the function $\sub()$ applied to a smaller problem instance.

\question[2]
Provide a recurrence for $\sub(n)$. Your answer to the previous part should be useful.
\question[5]
Using the recurrence, develop a memoized algorithm to compute $\sub(n)$. Remember that a memoized algorithm is always recursive. Your algorithm can use the quantities $\subsum(l,r)$ (since these values can be pre-computed by your algorithm from part 1 above).

\question[2]
What is the runtime of your memoized algorithm?  Provide a short justification.

\question[3]
Bonus: Suppose that instead of using the $\subsum()$ function to score a subarray, now use a linear variant:
\[
    \subsum'(l,r) = bx + c,
\]
where $b, c \in \mathbb{Z}$ and $x$ is the sum of all the elements in the sub-array $A[l..r]$. For this variant, we want to find the score of a partition that maximizes the sum of all subarrays.  Explain how we can achieve this with an algorithm whose runtime is faster than that for the original problem.

\end{questions}

\newpage

%=========================================================================================
\section{Legend of Zelda}

[Thanks to Denis Lalaj for contributing this problem.]

\vspace{.1in}

See the dynamic programming tutorial for motivation for this problem. An instance of the problem is $G[1..m][1..n]$---a 2D array of size $m \times n$ where each entry $G[i,j]$ is an integer (either positive or negative), representing the points  that Link can accumulate in the room at coordinates $(i,j)$. 

Link starts his quest at room $(1,1)$ and needs to get to the $(m,n)$ corner, via a path where each move is either to the right or downwards. Link starts with some initial number of points, and upon entering each room (including the first), Link's points are adjusted up or down by adding the points in that room. Link must ensure that he always has at least one point.

For $1 \le i \le m$ and $1\le j \le m$, let $HP[i,j]$ denote the minimum number of points that Link needs to have, when starting by entering the room with coordinates $(i, j)$, in order to reach $(m,n)$ while always having at least one point. The problem is to compute $HP[1,1]$---the minimum number of points needed initially in the full game.

In the tutorial, you studied the following recurrence for $HP[i,j]$:

\[
HP[i,j] = \left\{\begin{array}{ll}
\max(1, 1-G[m,n]), & \mbox{ if  $i=m$ and  $j=n$},\\
\max(1, \min(HP[i+1,j], HP[i,j+1]) -  G[i,j]), & \mbox{ if $1 \le i \le m-1$ or $1\le j \le n-1$}, \\
\infty,  & \mbox{ if $i = m+1$ or $j = n+1$}.
\end{array}
\right.
\]

\begin{questions}
\question[5]
Using this recurrence, 
develop a dynamic programming algorithm  to compute $HP(1,1)$. Remember that a dynamic programming algorithm is iterative (involving loops), not recursive.

\question[2]
What is the runtime of each of your algorithm?  Provide a short justification.

\end{questions}

\newpage

%=========================================================
\section{Subset Sums mod $m$}
In this example of dynamic programming, the recurrence describes a {\em set} of values, rather than a single value as in earlier examples.
An instance $I$ of the problem is a set $\{x_1, x_2, \ldots, x_n\}$ of nonnegative integers and a positive integer $m$. We say that a value $v$, with $0 \le v \le m-1$, is \emph{feasible} with respect to instance $I$ if for some non-empty subset $R$ of $\{x_1, x_2, \ldots, x_n\}$, 
\[ 
\sum_{x \in R} x \equiv v \pmod m.  
\]
Also, for each $i, 0\le i \le n$, let $V(i)$ be the set of feasible values with respect to instance $( \{x_1, x_2, \ldots, x_i\}, m)$.

\begin{questions}
\question[3]
Explain why the following recurrence holds for $V(i)$.
\[
   V(i) = \left\{
   \begin{array}{ll}
   \emptyset, & \mbox{ if } i = 0 \\
   V(i-1) \hspace{.1in} \cup \hspace{.1in} \cup_{v \in V(i-1)} \{ (v + x_i) \pmod m \} \hspace{.1in} \cup \hspace{.1in} \{ x_i \pmod m \}, & \mbox{ if } 1 \le i \le n.
   \end{array}
   \right.
   \]
   

\question[5]
\label{dp}
Design a dynamic programming algorithm that, given an instance $I = (\{x_1,x_2, \ldots, x_n\}, m)$ of Subset Sums mod $m$, determines whether 0 is feasible with respect to $I$. The algorithm outputs "Yes" or "No".  Your algorithm should run in $O(nm)$ time. Your pseudocode should create a solution array, Soln$[0..n]$, and should store the set of values $V(i)$ in Soln$[i]$. The pseudocode can use set operations such as $\cup$ as in the recurrence above.

\question[2]
Explain why your algorithm of part \ref{dp} runs in time $O(n m)$.

\question[3]
\label{reconstruct}
Suppose that indeed there is a subset $R$ of $S$ such that $\sum_{x  \in R} x \equiv 0 \pmod m$.  Fill in the missing parts (indicated by ????) of the following algorithm, that finds such a subset
$R$ by making a call to \textsc{Find-$R$($n,0$)}.
You can assume that the array Soln$[0..n]$ has already been pre-computed, where Soln$[i]$ stores the set $V(i)$; the algorithm uses the array Soln. Your algorithm should run in time $O(n)$.

\begin{algorithmic}
\Procedure{Find-$R$}{$i,v$}
   \State $\triangleright$ return non-empty $R \subseteq \{x_1,\ldots,x_i\}$ for which $\sum_{x \in R} x \equiv v \pmod m$, {\em given} that $R$ exists
   \State  $\triangleright$ the algorithm has access to the array Soln$[0..n]$
   \If{$x_i \pmod m = v$}
      \State  ????
    \ElsIf{$v \in \mbox{Soln}[i-1]$}
         \State \Return ????
    \Else
         \State $v' \gets v - x_i \pmod m$
         \State \Return $\{ x_i \} \hspace{.1in} \cup$ ????
    \EndIf
\EndProcedure
\end{algorithmic}

\question[2]
Explain why your algorithm of part \ref{reconstruct} runs in time $O(n)$.

\end{questions}


\end{document}